\documentclass[11pt,]{article}
\usepackage[margin=1in]{geometry}
\newcommand*{\authorfont}{\fontfamily{phv}\selectfont}
\usepackage[]{mathpazo}

\usepackage{amssymb,amsmath}
\usepackage{xcolor}

%TIKZ and Flow chart material
\usepackage{tikz}
\usetikzlibrary{shapes.geometric, arrows}
\tikzstyle{startstop} = [rectangle, rounded corners, minimum width=3cm, minimum height=1cm,text centered, draw=black, fill=red!30]
\tikzstyle{io} = [trapezium, trapezium left angle=70, trapezium right angle=110, minimum width=3cm, minimum height=1cm, text centered, draw=black,fill=white]
\tikzstyle{process} = [rectangle, minimum width=3cm, minimum height=1cm, text centered, draw=black, fill=white]
\tikzstyle{decision} = [diamond, minimum width=3cm, minimum height=1cm, text centered, draw=black, fill=white]
\tikzstyle{smalldecision} = [diamond, minimum width=1cm, minimum height=1cm, text centered, draw=black, fill=white]
\tikzstyle{arrow} = [thick,->,>=stealth]
\tikzstyle{line} = [thick,-]
\tikzstyle{dot} = [circle,inner sep=0.5pt,draw=black, fill=black]


%\tikzstyle{startstop} = [rectangle, rounded corners, minimum width=3cm, minimum height=1cm,text centered, %draw=black, fill=red!30]
%\tikzstyle{io} = [trapezium, trapezium left angle=70, trapezium right angle=110, minimum width=3cm, %minimum height=1cm, text centered, draw=black,fill=white]
%\tikzstyle{process} = [rectangle, minimum width=3cm, minimum height=1cm, text centered, draw=black, %fill=white]
%\tikzstyle{decision} = [diamond, minimum width=3cm, minimum height=1cm, text centered, draw=black, %fill=green!30]
%\tikzstyle{arrow} = [thick,->,>=stealth]




\usepackage{abstract}
\renewcommand{\abstractname}{}    % clear the title
\renewcommand{\absnamepos}{empty} % originally center
\newcommand{\blankline}{\quad\pagebreak[2]}

\providecommand{\tightlist}{%
  \setlength{\itemsep}{0pt}\setlength{\parskip}{0pt}}
\usepackage{longtable,booktabs}

\usepackage{parskip}
\usepackage{titlesec}
\titlespacing\section{0pt}{12pt plus 4pt minus 2pt}{6pt plus 2pt minus 2pt}
\titlespacing\subsection{0pt}{12pt plus 4pt minus 2pt}{6pt plus 2pt minus 2pt}

\titleformat*{\subsubsection}{\normalsize\itshape}

\usepackage{titling}
\setlength{\droptitle}{-.25cm}

%\setlength{\parindent}{0pt}
%\setlength{\parskip}{6pt plus 2pt minus 1pt}
%\setlength{\emergencystretch}{3em}  % prevent overfull lines

\usepackage[T1]{fontenc}
\usepackage[utf8]{inputenc}

\usepackage{fancyhdr}
\pagestyle{fancy}
\usepackage{lastpage}
\renewcommand{\headrulewidth}{0.3pt}
\renewcommand{\footrulewidth}{0.0pt}
%\lhead{\footnotesize Problem Set \#3 }
\lhead{\footnotesize BUEC 311: Problem Set \#3, Shifting equilibria}
\rhead{\footnotesize \today}
\lfoot{\small \copyright }
\cfoot{}
\rfoot{\small \thepage/\pageref*{LastPage}}

\fancypagestyle{firststyle}
{
\renewcommand{\headrulewidth}{0pt}%
   \fancyhf{}
   \fancyfoot[C]{\thepage/\pageref*{LastPage}}
}

%\def\labelitemi{--}
%\usepackage{enumitem}
%\setitemize[0]{leftmargin=25pt}
%\setenumerate[0]{leftmargin=25pt}




\makeatletter
\@ifpackageloaded{hyperref}{}{%
\ifxetex
  \usepackage[setpagesize=false, % page size defined by xetex
              unicode=false, % unicode breaks when used with xetex
              xetex]{hyperref}
\else
  \usepackage[unicode=true]{hyperref}
\fi
}
\@ifpackageloaded{color}{
    \PassOptionsToPackage{usenames,dvipsnames}{color}
}{%
    \usepackage[usenames,dvipsnames]{color}
}
\makeatother
\hypersetup{breaklinks=true,
            bookmarks=true,
            pdfauthor={},
             pdfkeywords = {},
            pdftitle={Problem Set \#3: Shifting equilibria},
            colorlinks=true,
            citecolor=blue,
            urlcolor=blue,
            linkcolor=magenta,
            pdfborder={0 0 0}}
\urlstyle{same}  % don't use monospace font for urls


\setcounter{secnumdepth}{0}


%



\usepackage{setspace}

\title{\vspace{-1.5cm}\Large{BUEC 311: Business Economics, Organization
and Management}\medskip\\\Large{Problem Set \#3}
\medskip\\\Large{Shifting equilibria}
}
\date{\vspace{-.75cm}\Large{\today}}

\definecolor{light-gray}{gray}{0.8}


\begin{document}

\vspace{-5cm}\maketitle
 \tikz [remember picture,overlay]
    %\node[yshift=-1.65cm,xshift=0cm] at (current page.north)
    %\node[yshift=-1.65cm,xshift=0cm] at (current page.north)
        %or: (current page.center)
        \node[yshift=-1cm,xshift=6.5cm] at (current page.north west)
        %{\includegraphics[width=3in]{UA-ASB-COLOUR.png}};
        {\includegraphics[width=.5\paperwidth]{../images/UA-ASB-COLOUR.png}};
\vspace{-.75cm}		
		\thispagestyle{firststyle}



This week's problems are meant to give you more practice on the basic
tricks you'll need to manipulate supply and demand curves and calculate
equilibrium conditions, but these also involve shifts in equilibria. Be
methodical about your calculations and your analysis. In these
questions, I've again given you multiple steps intentionally to build
habits like checks and double-checks that you'll want to use. I won't
always give these to you on exam questions.

\begin{enumerate}
\def\labelenumi{\arabic{enumi}.}
\tightlist
\item
  Let's use the market demand curve for coffee we derived in the last
  problem set of \(Q=4000-\frac{2000p}{3}\). Let's also assume that the
  marginal cost of supplying coffee to the market is
  \(Q=\frac{2000\,p}{3}\).

  \begin{enumerate}
  \def\labelenumii{\alph{enumii})}
  \tightlist
  \item
    Calculate the equilibrium price and quantity (a repeat from last
    week)
  \item
    Suppose that a drought reduces the coffee crop increasing the cost
    of supplying coffee leading to a reduction in the supply function to
    \(Q=\frac{1000\,p}{3}\). Solve for the new equilibrium price and
    quantity in the market
  \item
    Show, on a graph, the changes in the inverse supply and demand
    functions, and label the original and new equilibrium prices and
    quantities.
  \item
    Is this a change in demand or a movement along the demand curve?
    Why?
  \end{enumerate}
\end{enumerate}

\newpage

\begin{enumerate}
\def\labelenumi{\arabic{enumi}.}
\setcounter{enumi}{1}
\tightlist
\item
  Use the market demand curve for coffee we derived in the last problem
  set of \(Q=4000-\frac{2000p}{3}\) along with the supply function of
  \(Q=\frac{2000\,p}{3}\).

  \begin{enumerate}
  \def\labelenumii{\alph{enumii})}
  \tightlist
  \item
    Suppose a coffee tax of \$0.50 per cup is levied by the university
    in an attempt to recover some revenues lost to COVID-19. What's the
    impact on the market if you assume that the tax is charged to the
    coffee vendor based on the number of cups of coffee sold?
  \item
    What happens if the same tax is added at the cash register and
    charged to patrons directly, with the proceeds of the tax remitted
    to the university?
  \item
    Why does the equilibrium price of coffee not increase by the same
    amount as the tax?
  \end{enumerate}
\end{enumerate}

\newpage

\begin{enumerate}
\def\labelenumi{\arabic{enumi}.}
\setcounter{enumi}{2}
\tightlist
\item
  On slide 81 of the full supply and demand deck, I give you a graph of
  a concurrent shift in supply and demand due to a change in crude oil
  prices affecting both the input costs to produce gasoline and the
  income of consumers. Based on the remainder of the information in the
  deck:

  \begin{enumerate}
  \def\labelenumii{\alph{enumii})}
  \tightlist
  \item
    Write out the initial supply and demand functions and solve for the
    initial equilibrium price and quantity (this is a freebie, as this
    is all in the deck)
  \item
    Solve for the price and quantity at the new equilibrium (not in the
    deck)
  \item
    Now, assume the government levies a 20c/l tax (\$0.20/l in the units
    in the deck) on gasoline payable by suppliers. Calculate a new
    equlibrium price and quantity.
  \item
    Has the price increased by the cost of the tax? If not, why did the
    producers not simply \textit{pass the cost on} to consumers?
  \item
    Now, assume the government levies a 20c/l tax (\$0.20/l in the units
    in the deck) on gasoline payable at the cash register by consumers.
    Calculate the new equlibrium price and quantity.
  \item
    In the new equilibrium, what's the revenue per litre earned by
    producers and what is the price paid by consumers?\\
    \newpage
  \end{enumerate}
\item
  In your own words and graphs (using the supply and demand functions
  from the class deck), describe the impact of a price ceiling on
  gasoline applied at a cost of \$0.80/l. What is likely to result from
  this price ceiling? (For now, assume the market is competitive. We'll
  abandon this in due course.)
\end{enumerate}




\end{document}
